\chapter{Оптимизация кода}\label{ch:optimization}

Одной из ключевых задач при разработке современного веб-приложения является оптимизация объема передаваемых данных и времени загрузки.
В рамках данного проекта была реализована поэтапная стратегия снижения веса приложения и улучшения пользовательского опыта.

\section{Code Splitting с использованием React Router DOM}\label{sec:code-splitting}

Первым этапом стала реализация динамической загрузки модулей (code splitting), основанная на возможностях библиотеки React Router DOM\@.

%Вместо загрузки всего приложения сразу, каждый маршрут стал загружаться только при необходимости — при переходе пользователя к соответствующему разделу.
Вместо загрузки всего приложения сразу сначала загружаются только страницы авторизации и регистрации врача, а дальше при успешной авторизации догружаются нужные страницы для двух ролей: врача и пациента.
Это позволило существенно сократить первоначальный объем загружаемого кода.

До внедрения механизма отложенной загрузки общий размер приложения составлял около 3 МБ.
После внедрения code splitting и разделения кода по маршрутам, вес основных бандлов снизился до 1.9 МБ, что значительно улучшило производительность при первом запуске.

\section{Сжатие Gzip}\label{sec:gzip}

Следующим этапом стала настройка сжатия данных на стороне сервера с использованием алгоритма Gzip.
Сжатие применялось ко всем статическим ресурсам, включая JavaScript, CSS и JSON-файлы.

Результаты показали, что в среднем объём передаваемых данных снизился примерно на 70\%: вес основного бандла составил около 400 КБ вместо 1.9 МБ.
Это существенно ускорило загрузку интерфейса даже при медленном интернет-соединении.

\section{Кэширование и контроль версий}\label{sec:caching}

Также была реализована стратегия кэширования: статические файлы запрашиваются браузером повторно только при их фактическом обновлении на сервере (например, после новой сборки). Это стало возможным благодаря использованию уникальных хешей в именах файлов (например, `main.abc123.js`), что гарантирует их переиспользование до следующей сборки приложения.

Таким образом, приложение стало не только легче, но и «умнее» в плане сетевой активности: ресурсы повторно не скачиваются без необходимости, а сами запросы к API также сократились по размеру за счёт минимизации тела ответов и заголовков.


В результате проделанных шагов приложение стало значительно быстрее и эффективнее с точки зрения сетевого трафика и загрузки интерфейса, особенно на мобильных устройствах и в условиях нестабильного соединения.
